\section{Conditional forest bounds in large clean blocks}
\label{sec:path-high-girth-clean-blocks}

The last step of the forest argument can use two conditional probability bounds instead of an inverse second moment.
This requires only that a block's girth be much larger than the logarithm of its boundary bound.
All logarithms have base two.

\begin{lemma}[A conditional bound with rank larger than the port scale]
\label{lem:path-clean-core-affine-forest}
There is an absolute constant $C_0$ with the following property.
Let $L\ge4$, and let $N$ be a connected normalized $k$-pole of rank $h$ such that
\[
 h\ge L,\qquad k\le2L,
 \qquad\operatorname{girth}(N)\ge D\ge C_0\log(L+2).
\]
Then every even port subset $T$ has affine forest proportion at most $3/5$.
The assertion also permits $k=0,1$, when the only even port subset is empty.
\end{lemma}

\begin{proof}
If $k\le1$, the only relevant class is the even space.
Its only forest is the empty state, of proportion $2^{-h}\le1/16$.
Assume $k\ge2$.
Since $h\ge L\ge4$ and $k\le2L\le2h$, the stationary branching entropy satisfies
\[
 \sigma=\frac{3(h-1)}{3(h-1)+k}\ge\frac12.
\]
The entropy Moore bound and $n=2h+k-2\le4h$ give
\[
 D\le4+4\log n\le12+4\log h,
 \qquad h\ge2^{D/4-3}.
\]

Put $s=\lfloor D/8\rfloor$.
The radius-$s$ balls about all ports contain at most $3k2^s\le6L2^{D/8}$ vertices.
For a sufficiently large absolute $C_0$, the hypothesis on $D$ makes this less than $n\ge h\ge2^{D/4-3}$.
There is therefore a radius-$s$ ball avoiding every port.
Since $2s<D$, its cubic breadth-first expansion has at least $2^s\ge2^{D/8-1}$ vertices.
Relative to any fixed root, this ball meets at most $2s+1\le D$ distance layers.
One layer consequently has at least
\[
 \frac{2^{D/8-1}}D
\]
vertices.

If an affine forest proportion were greater than $3/5$, the BFS marker and Lemma~\ref{lem:path-thin-marker-layer-width}, with $\varepsilon=1/10$, would bound every layer by
\[
 K\left[100(\log40)^{11}+(k+1)^2\right]
 \le K_2(L+2)^2
\]
for an absolute $K_2$.
Increasing $C_0$ makes the preceding lower bound larger than this upper bound whenever $D\ge C_0\log(L+2)$.
To verify uniformity, $2^{D/8}/D$ is increasing for $D\ge16$, and at $D=C_0\log(L+2)$ its ratio to $(L+2)^2$ is
\[
 \frac{(L+2)^{C_0/8-2}}{C_0\log(L+2)}.
\]
This can be made arbitrarily large uniformly for $L\ge4$ by increasing $C_0$.
The contradiction proves the lemma.
\end{proof}

\begin{lemma}[A large block contains a suitable core]
\label{lem:path-clean-block-core}
Let $G$ be connected and cubic, and let $R$ be a nonempty proper vertex set with
\[
 |R|\ge8L^2,
 \qquad |\delta_G R|\le2L,
 \qquad L\ge4.
\]
Then a connected component of $G[R]$ has a connected two-core of rank at least $L$ and at most $2L$ degree-two vertices.
Its other vertices have degree three.
\end{lemma}

\begin{proof}
Write $n=|R|$, $d=|\delta_G R|$, and $c=\operatorname{comp}(G[R])$.
Every component has a boundary edge, so $1\le c\le d\le2L$.
The degree sum gives
\[
 \beta(G[R])=\frac{n-d}{2}+c.
\]
The average component rank is therefore at least
\[
 \frac{n-d}{2d}+1
 \ge\frac{8L^2-2L}{4L}+1>L.
\]
Choose a component of rank at least $L$.
Pruning its degree-one vertices preserves rank and leaves a nonempty connected two-core.
The core has minimum degree two and maximum degree three.

At each pruning step, a removed leaf has two incidences outside the current core candidate and one incidence inside it.
Its removal decreases the boundary by one.
Thus the final core has boundary at most the original component boundary, which is at most $2L$.
Every degree-two core vertex has exactly one boundary incidence, since $G$ is cubic; degree-three vertices have none.
The asserted port bound follows.
\end{proof}

\begin{proposition}[Two clean blocks force probability below one half]
\label{prop:path-two-clean-blocks}
Let $G$ be connected and cubic, and let $P$ be any marked set.
Suppose $G-P$ contains two vertex-disjoint sets $R_1,R_2$ such that, for $i=1,2$,
\[
 |R_i|\ge8L^2,
 \qquad |\delta_G R_i|\le2L,
 \qquad\operatorname{girth}(G[R_i])\ge D,
\]
where $L\ge4$ and $D\ge C_0\log(L+2)$.
Then
\[
 \Pr(\xi_G-P\text{ is a forest})\le\frac9{25}<\frac12.
\]
In particular, these hypotheses hold eventually for any sequence with $D/\log(L+2)\to\infty$.
\end{proposition}

\begin{proof}
Choose cores $K_i$ by Lemma~\ref{lem:path-clean-block-core}.
They are vertex-disjoint, avoid $P$, and retain girth at least $D$.
Condition on every edge outside $E(K_1)\cup E(K_2)$.
For every conditioning of positive probability, the two remaining edge sets are independent uniform members of their respective affine parity classes.
Indeed, each boundary parity is fixed, and the solution set is the direct product of the two affine internal cycle spaces.
Only degree-two core vertices can have odd prescribed parity.
Their odd sets have even cardinality.

Lemma~\ref{lem:path-clean-core-affine-forest} bounds each conditional forest probability by $3/5$.
The forest event outside $P$ implies that both core restrictions are forests.
Its conditional probability is therefore at most $(3/5)^2$.
Averaging over the conditioned edges gives the result.
Edges between the two cores cause no difficulty, since they are among the conditioned edges.
\end{proof}
